\section{Parameterized speed Profile}
\label{sec:speed}
% 本研究构建面向单次锻造的锻压速度曲线决策器，能够根据锻件初始温度场快速生成完整速度曲线。参考前人工作\cite{lin2011critical}，将速度曲线限定为匀速-加速-匀速模式，并采用四参数平滑函数进行参数化，其数学表达式为：
% \begin{equation}
% v_{y}(t)=\begin{cases}
% v_{0}, & t<t_{0} \\
% v_{0}+\left(v_{1}-v_{0}\right)\xi^{3}\left(10-15\xi+6\xi^{2}\right), & t\geq t_{0}
% \end{cases}
% \end{equation}
% 其中$\xi=(t-t_0)/(t_1-t_0)$。决策变量为 $\theta=\{v_0,v_1,t_0,t_1\}$，分别表示初始速度、目标速度、加速起止时间（图\ref{fig:speed_profile}）。

% \begin{figure}[htbp]
% \centerline{\includegraphics[width=\linewidth]{speed_profile.eps}}
% \caption{Parameterized speed profile.}
% \label{fig:speed_profile}
% \end{figure}

% 该函数在$t_0$和$t_1$处具有二阶连续导数，保证加速过程平滑。参数化实现了从低维参数到连续曲线的映射$v_y (t;\theta)$，将速度优化问题转化为对$\theta$的搜索问题。
This study constructs a forging-speed profile decision system for single
forging operations, capable of rapidly generating a complete speed
profile based on the initial temperature field of the workpiece. In line
with previous work \cite{lin2011critical}, the speed profile is defined
by a constant-acceleration-constant pattern and parameterized using a
smooth four-parameter function. Its mathematical expression is given by:
\begin{equation}
v_{y}(t)=\begin{cases}
v_{0}, & t<t_{0} \\
v_{0}+\left(v_{1}-v_{0}\right)\xi^{3}\left(10-15\xi+6\xi^{2}\right), & t\geq t_{0}
\end{cases}
\end{equation}
where $\xi=(t-t_0)/(t_1-t_0)$. The decision variables are
$\theta=\{v_0, v_1, t_0, t_1\}$, denoting the initial speed, target
speed, and the start and end times of the acceleration phase,
respectively (Fig. \ref{fig:speed_profile}).

\begin{figure}[htbp]
\centerline{\includegraphics[width=\linewidth]{velocity_profile.eps}}
\caption{Parameterized speed profile.}
\label{fig:speed_profile}
\end{figure}

The function possesses a continuous second-order derivative at both
$t_0$ and $t_1$, guaranteeing a smooth acceleration process. This
parameterization effectively maps the low-dimensional parameter set
$\theta$ to a continuous curve, $v_y (t;\theta)$, thereby transforming
the speed profile optimization problem into a search for the optimal
$\theta$.
